\section{Continue from filtered bytes to reconstructed samples}
The inflater now handles all three block kinds. Its output still contains PNG filter
selectors and residuals. This next construction changes that buffer in place, following
the row interface above. RGB packing will receive these reconstructed samples later.
The two operations remain separate: no display word is written during reconstruction.

Attach \texttt{drafts/png-row-construction/attach.py} to your completed Huffman-stage
program. It preserves your earlier instruction edits and extra functions, replaces
only the supplied temporary wrapper, and appends thirteen unfinished regions. The
new wrapper calls scalar or MSA reconstruction according to the existing mode control.
It still stops before packing and returns development status after valid rows.
Generate the selected fixture with \texttt{fixtures.py}; use \texttt{prepare.py} with
the recipe ID to select its path and mode without replacing your instructions.

The zero upper row, row cursors and predictor temporaries below describe the supplied
construction layout. They make individual observations reproducible. A free-writing
implementation may choose other registers, stack offsets or predictor formulations,
provided it meets the callable interface, access bounds and failure guarantees.
An unfinished region returns development status and releases an allocated working
frame; calling a framed internal label directly would bypass that ownership rule.
